\chapter{Particle Mechanics, Four-Momentum, and Rest Energy}

\section{The Approximation We Will Use Repeatedly}

The mechanics of a relativistic particle must answer familiar questions in a new language. How does a particle move? What are its momentum and energy? Why do those quantities fit together into a four-vector, and how does \(E=mc^2\) emerge from the ordinary machinery of mechanics? Before taking up those questions, we need one elementary approximation that will recur throughout the calculation.

For any exponent \(p\), the binomial series is
\begin{equation}
\begin{aligned}
(1+\epsilon)^p
&=1+p\epsilon\\
&\quad+\frac{p(p-1)}{2!}\epsilon^2\\
&\quad+\frac{p(p-1)(p-2)}{3!}\epsilon^3
+\cdots .
\end{aligned}
\label{eq:sr3-binomial}
\end{equation}
The exponent may be positive, negative, fractional, or even complex. For sufficiently small \(\epsilon\), the series converges rapidly and each successive term is smaller. We may then stop at whatever order the desired accuracy requires. The zeroth-order term is often the largest, but in the applications below it is frequently a constant and therefore dynamically uninteresting. The first correction is where the motion appears.

\begin{figure}[htbp]
\centering
\includegraphics[width=0.90\textwidth]{lecture_03_figure_01.png}
\caption{The binomial series through cubic order, used to recover nonrelativistic mechanics from relativistic formulas.}
\lectureframe{00:04:24}
\end{figure}

We normally set \(c=1\), so a velocity \(v\) means the fraction \(v/c\) of light speed. Take \(\epsilon=-v^2\). With \(p=\tfrac12\),
\begin{equation}
\sqrt{1-v^2}
=
1-\frac{v^2}{2}-\frac{v^4}{8}-\cdots
\approx
1-\frac{v^2}{2},
\qquad |v|\ll 1.
\label{eq:sr3-sqrt-expansion}
\end{equation}
With \(p=-\tfrac12\), the two minus signs cancel:
\begin{equation}
\frac{1}{\sqrt{1-v^2}}
=
1+\frac{v^2}{2}+\frac{3v^4}{8}+\cdots
\approx
1+\frac{v^2}{2}.
\label{eq:sr3-gamma-expansion}
\end{equation}
Many factors of \(1/2\) in classical mechanics can be recognized, after the fact, as the first correction in one of these relativistic expansions. The approximation fails as \(v\) approaches one; then no finite truncation captures the singular behavior of the exact expression.

\begin{figure}[htbp]
\centering
\includegraphics[width=0.90\textwidth]{lecture_03_figure_02.png}
\caption{Both low-velocity formulas are visible: \(\sqrt{1-v^2}\approx1-v^2/2\) and \((1-v^2)^{-1/2}\approx1+v^2/2\).}
\lectureframe{00:07:39}
\end{figure}

\section{What Counts as a Particle?}

A particle in mechanics need not be elementary. It may be an electron, the Sun, a box of chocolates, or a donut. What matters is that the object holds together well enough for its overall motion to be described by a center of mass. A donut is a useful reminder that the center of mass need not lie in the material itself.

In Newtonian mechanics the center-of-mass variables obey
\begin{equation}
\vec p=m\vec v,
\qquad
K=\frac12mv^2,
\qquad
\vec F=m\vec a.
\label{eq:sr3-newtonian-particle}
\end{equation}
These formulas do not care whether the object is elementary or composite. A composite object also contains internal energy. A moving donut, for example, has center-of-mass kinetic energy plus chemical, thermal, and binding energy:
\begin{equation}
E_{\text{Newtonian}}
=
E_{\text{internal}}+\frac12mv^2.
\label{eq:sr3-newtonian-internal}
\end{equation}
The internal term is independent of the overall center-of-mass velocity. Relativity will retain this distinction but will reveal that the internal energy contributes to the invariant mass of the whole object.

The tools required are the same ones used in classical mechanics: trajectories, the principle of stationary action, Lagrangians, canonical momenta, Hamiltonians, and the relation between continuous symmetries and conservation laws.

\section{Worldlines and Causal Separation}

Plot a system's position horizontally and time vertically. Its history is then a curve in spacetime. For one particle this curve is its \emph{worldline}. The concept exists in Newtonian mechanics, but it becomes fundamental in relativity because Lorentz transformations mix space and time.

Consider two events separated by
\begin{equation}
\Delta x^\mu=(\Delta t,\Delta\vec x),
\qquad
\Delta\vec x^{\,2}
=
(\Delta x)^2+(\Delta y)^2+(\Delta z)^2.
\end{equation}
In units \(c=1\), their invariant interval is
\begin{equation}
(\Delta\tau)^2
=
(\Delta t)^2-\Delta\vec x^{\,2}.
\label{eq:sr3-interval}
\end{equation}

\subsection{Timelike separation}

If \((\Delta t)^2>\Delta\vec x^{\,2}\), then \((\Delta\tau)^2>0\) and the separation is timelike. There is a Lorentz frame in which the two events occur at the same spatial point. In that frame the line joining them is vertical and a clock can travel from one event to the other. Its elapsed reading is the proper time \(\Delta\tau\).

\begin{figure}[htbp]
\centering
\includegraphics[width=0.90\textwidth]{lecture_03_figure_03.png}
\caption{A timelike interval has more temporal than spatial separation; a rest frame exists in which the two events share one position.}
\lectureframe{00:14:42}
\end{figure}

\subsection{Spacelike separation}

If \(\Delta\vec x^{\,2}>(\Delta t)^2\), the interval is spacelike. No inertial frame places both events at the same spatial point, but there is a frame in which they are simultaneous. Other frames can reverse their time order. Therefore neither \emph{event \(A\) happened first} nor \emph{event \(B\) happened first} is an invariant statement for a spacelike pair. This is the relativity of simultaneity in geometric form.

\begin{figure}[htbp]
\centering
\includegraphics[width=0.90\textwidth]{lecture_03_figure_04.png}
\caption{For a spacelike interval, a change of inertial frame can make the two events simultaneous or reverse their temporal order.}
\lectureframe{00:17:01}
\end{figure}

A massive particle has a timelike worldline: every sufficiently short segment satisfies \(d\tau^2>0\). Successively composing velocities of \(0.9c\) may bring a particle ever closer to \(c\), but never beyond it. The approach resembles Zeno's paradox. This causal restriction applies not only to material objects but also to signals and information, because information must be carried by physical systems. A faster-than-light signal would allow different frames to disagree about which end of its history is the cause and which is the effect.

\section{Four-Velocity}

Focus on a short timelike segment. In conventional units,
\begin{equation}
d\tau
=
\sqrt{dt^2-\frac{d\vec x^{\,2}}{c^2}}.
\label{eq:sr3-proper-time-dimensional}
\end{equation}
The factor of \(c^2\) is forced by dimensions: \(dt^2\) has units of time squared, whereas \(d\vec x^{\,2}\) has units of length squared.

\begin{classroomqa}
\audiencequestion{Where must the speed of light be inserted when the interval is written in ordinary units?}
\lecturerresponse{The spatial term must be divided by \(c^2\), giving \(d\tau^2=dt^2-d\vec x^{\,2}/c^2\). We will usually return to \(c=1\), but dimensional consistency always recovers the missing factors. \lecturetimestamp{00:21:23}}
\end{classroomqa}

\begin{figure}[htbp]
\centering
\includegraphics[width=0.90\textwidth]{lecture_03_figure_05.png}
\caption{A worldline displacement \(\Delta x^\mu=(\Delta t,\Delta\vec x)\) and the proper-time increment associated with it.}
\lectureframe{00:21:03}
\end{figure}

Ordinary velocity has three spatial components,
\begin{equation}
v^i=\frac{dx^i}{dt},
\qquad i=1,2,3.
\end{equation}
The relativistic velocity is obtained by dividing the full spacetime displacement by invariant proper time:
\begin{equation}
u^\mu
=
\frac{dx^\mu}{d\tau},
\qquad
\mu=0,1,2,3.
\label{eq:sr3-four-velocity-definition}
\end{equation}
Greek indices run over spacetime; Latin indices run over space. The zeroth component of position is \(x^0=t\) in \(c=1\) units, so
\begin{equation}
u^0=\frac{dt}{d\tau}.
\end{equation}

\begin{figure}[htbp]
\centering
\includegraphics[width=0.90\textwidth]{lecture_03_figure_06.png}
\caption{Ordinary velocity divides spatial displacement by coordinate time, while four-velocity divides all four components by proper time.}
\lectureframe{00:24:05}
\end{figure}

To relate \(u^\mu\) to \(\vec v\), factor \(dt^2\) from the interval:
\begin{align}
d\tau^2
&=dt^2-d\vec x^{\,2}\\
&=dt^2\left(1-\frac{d\vec x^{\,2}}{dt^2}\right)
=dt^2(1-v^2).
\end{align}
Taking the positive root for a future-directed worldline,
\begin{equation}
\frac{d\tau}{dt}=\sqrt{1-v^2},
\qquad
\frac{dt}{d\tau}=\frac{1}{\sqrt{1-v^2}}
\equiv\gamma.
\label{eq:sr3-time-component}
\end{equation}
For a spatial component,
\begin{equation}
v^i
=
\frac{dx^i}{dt}
=
\frac{dx^i}{d\tau}\frac{d\tau}{dt}
=u^i\sqrt{1-v^2}.
\end{equation}
Thus
\begin{equation}
\boxed{
u^\mu
=
\gamma(1,\vec v)
},
\qquad
u^0=\gamma,
\qquad
u^i=\gamma v^i.
\label{eq:sr3-four-velocity-components}
\end{equation}
At low speed \(u^0\to1\) and \(u^i\to v^i\). Near light speed both \(u^0\) and the magnitude of the spatial four-velocity become large even though \(|\vec v|<1\).

\begin{figure}[htbp]
\centering
\includegraphics[width=0.90\textwidth]{lecture_03_figure_07.png}
\caption{The settled relation between ordinary velocity and all four components of \(u^\mu\).}
\lectureframe{00:30:40}
\end{figure}

\subsection{Why only three components are independent}

Four-velocity has four components but only three independent degrees of freedom. Dividing the interval identity by \(d\tau^2\) gives
\begin{equation}
1
=
\left(\frac{dt}{d\tau}\right)^2
-\left(\frac{d\vec x}{d\tau}\right)^2,
\end{equation}
or
\begin{equation}
\boxed{
(u^0)^2-\vec u^{\,2}=1
}.
\label{eq:sr3-four-velocity-normalization}
\end{equation}
The one normalization condition removes one independent component.

\begin{figure}[htbp]
\centering
\includegraphics[width=0.90\textwidth]{lecture_03_figure_08.png}
\caption{The invariant interval immediately yields the normalization \((u^0)^2-\vec u^{\,2}=1\).}
\lectureframe{00:33:50}
\end{figure}

\begin{classroomqa}
\audiencequestion{Can the loss of one independent component be understood from differentiating with respect to proper time itself?}
\lecturerresponse{Yes. Proper time is not an independent fifth coordinate; it is fixed by the spacetime interval along the worldline. Dividing that defining relation by \(d\tau^2\) constrains the four derivatives to lie on the unit timelike hyperboloid \((u^0)^2-\vec u^{\,2}=1\). \lecturetimestamp{00:34:07}}
\end{classroomqa}

\section{An Invariant Action for a Free Particle}

What law governs a force-free worldline? We retain two lessons from classical mechanics:
\begin{enumerate}
\item the physical path makes an action stationary while its endpoints are held fixed;
\item the action is built additively from infinitesimal segments of the path.
\end{enumerate}
Relativity adds a third demand: every inertial observer must agree on the action. The natural infinitesimal scalar available on a timelike worldline is \(d\tau\). Therefore the free-particle action is
\begin{equation}
\boxed{
S=-m\int_a^b d\tau
}
\qquad(c=1).
\label{eq:sr3-invariant-action}
\end{equation}
The finite-sum picture is
\begin{equation}
S=-m\lim_{N\to\infty}\sum_{k=1}^{N}\Delta\tau_k.
\end{equation}
Since every observer agrees on each proper-time increment, every observer agrees on which path makes their sum stationary.

\begin{figure}[htbp]
\centering
\includegraphics[width=0.90\textwidth]{lecture_03_figure_09.png}
\caption{The action is built as a sum of invariant proper-time increments along a varied worldline with fixed endpoints.}
\lectureframe{00:38:30}
\end{figure}

\begin{classroomqa}
\audiencequestion{What happens to energy if the action is multiplied by an arbitrary constant?}
\lecturerresponse{Energy is not the starting invariant; it will be derived from the Lagrangian and Hamiltonian. Multiplying the entire action by a nonzero constant leaves its stationary paths unchanged, though it rescales the canonical quantities defined from that action. The normalization \(-m\) is chosen so those derived quantities match established Newtonian conventions. \lecturetimestamp{00:39:47}}
\end{classroomqa}

\begin{classroomqa}
\audiencequestion{Why is an invariant such as proper time a good candidate from which to build the action?}
\lecturerresponse{Because a stationary-action law should select the same physical path in every inertial frame. This is analogous to finding the shortest curve in a Euclidean plane: rotations change coordinate components but not each infinitesimal length, so they cannot change which curve is shortest. Replacing Euclidean length by proper time gives the relativistic construction. \lecturetimestamp{00:40:39}}
\end{classroomqa}

\begin{classroomqa}
\audiencequestion{Why choose the positive square root for proper time rather than the negative root?}
\lecturerresponse{Proper time is defined as the positive elapsed reading of a clock carried from past to future. Choosing the negative root consistently would merely reverse the sign of the action, changing a minimum into a maximum without changing the stationary worldline. In the clock's rest frame, the positive convention gives \(d\tau=dt\). \lecturetimestamp{00:42:59}}
\end{classroomqa}

The factor \(-m\) is likewise a normalization convention selected to recover Newtonian mechanics. Restoring units, the conventional action is
\begin{equation}
S=-mc^2\int_a^b d\tau.
\label{eq:sr3-action-dimensional}
\end{equation}

\begin{figure}[htbp]
\centering
\includegraphics[width=0.90\textwidth]{lecture_03_figure_10.png}
\caption{The finite proper-time sum becomes the invariant integral \(-m\int d\tau\) in natural units.}
\lectureframe{00:46:23}
\end{figure}

\section{The Relativistic Lagrangian}

To put the action into standard mechanical form, factor \(dt\) from \(d\tau\):
\begin{equation}
d\tau
=dt\sqrt{1-\dot{\vec x}^{\,2}},
\qquad
\dot{\vec x}^{\,2}
=\dot x^2+\dot y^2+\dot z^2.
\end{equation}
Then
\begin{equation}
S
=
\int_a^b L\,dt,
\qquad
\boxed{
L=-m\sqrt{1-v^2}
}
\quad(c=1).
\label{eq:sr3-relativistic-lagrangian}
\end{equation}
In conventional units,
\begin{equation}
L=-mc^2\sqrt{1-\frac{v^2}{c^2}}.
\label{eq:sr3-relativistic-lagrangian-c}
\end{equation}
This is the Lagrangian of a free particle: it depends on velocity but not position. No force or potential has yet been introduced.

\begin{figure}[htbp]
\centering
\includegraphics[width=0.90\textwidth]{lecture_03_figure_11.png}
\caption{Factoring out \(dt\) reveals the free relativistic Lagrangian as a function of the three velocity components.}
\lectureframe{00:49:57}
\end{figure}

\subsection{Recovering Lagrange and Newton}

Insert \eqref{eq:sr3-sqrt-expansion}:
\begin{align}
L
&=-m\sqrt{1-v^2}\\
&\approx-m\left(1-\frac{v^2}{2}\right)\\
&=-m+\frac12mv^2.
\label{eq:sr3-lagrangian-low-speed}
\end{align}
The constant \(-m\) is independent of both position and velocity. It contributes the same amount to every path with the same endpoint times and therefore drops out of the Euler--Lagrange equations. The part with dynamical content is precisely
\begin{equation}
L_{\mathrm{Newton}}=\frac12mv^2.
\end{equation}
Thus the sign and mass normalization in the invariant action were chosen correctly.

\begin{figure}[htbp]
\centering
\includegraphics[width=0.90\textwidth]{lecture_03_figure_12.png}
\caption{The low-speed expansion leaves a constant \(-m\) and the Newtonian free-particle Lagrangian \(mv^2/2\).}
\lectureframe{00:53:24}
\end{figure}

For the present free particle, the Euler--Lagrange equations will say only that the velocity is constant and the spacetime worldline is straight. The more revealing task is to identify the conserved momentum and energy.

\section{Canonical Momentum}

Because \(L\) has no position dependence, it is invariant under spatial translations. The corresponding conserved quantity is momentum. We do not assume the Newtonian formula \(p_i=mv_i\); we use the canonical definition
\begin{equation}
p_i=\frac{\partial L}{\partial\dot x^i}.
\label{eq:sr3-canonical-momentum}
\end{equation}
For \(L=-m\sqrt{1-v^2}\),
\begin{align}
p_i
&=
-m\frac{1}{2\sqrt{1-v^2}}
\frac{\partial}{\partial\dot x^i}(1-v^2)\\
&=
-m\frac{1}{2\sqrt{1-v^2}}(-2\dot x^i)\\
&=
\frac{mv_i}{\sqrt{1-v^2}}
=m\gamma v_i.
\label{eq:sr3-relativistic-momentum}
\end{align}
Since \(u^i=\gamma v^i\),
\begin{equation}
\boxed{
p^i=mu^i
}.
\label{eq:sr3-momentum-spatial}
\end{equation}
At low speed this reduces to \(p^i\approx mv^i\). Near light speed the denominator matters decisively.

\begin{figure}[htbp]
\centering
\includegraphics[width=0.90\textwidth]{lecture_03_figure_13.png}
\caption{Canonical differentiation of the square-root Lagrangian gives \(p_x=mv_x/\sqrt{1-v^2}\).}
\lectureframe{00:58:49}
\end{figure}

\subsection{Mass, not relativistic mass}

Here \(m\) is simply the particle's mass: an invariant label that does not change when the particle moves. Older texts sometimes call it \emph{rest mass} and call \(E/c^2\) a velocity-dependent \emph{relativistic mass}. Modern usage avoids that split. Motion changes energy and momentum, not invariant mass. The mass of an electron is the same parameter in every frame.

\section{From Three-Momentum to Four-Momentum}

A statement about one component of a four-vector is generally frame-dependent. For example, \(\Delta\vec x=0\) says that two events share a position in one frame, but a boost tilts that interval and makes \(\Delta\vec x'\ne0\). By contrast,
\begin{equation}
\Delta x^\mu=0
\end{equation}
is invariant: if all four components vanish in one frame, all four vanish in every frame. Likewise, \(\vec u=0\) only identifies the particle's rest frame; it is not an invariant description of the four-velocity.

For a closed system, ordinary momentum conservation reads
\begin{equation}
\vec P_{\mathrm{initial}}-\vec P_{\mathrm{final}}=0.
\label{eq:sr3-three-momentum-conservation}
\end{equation}
This is a three-vector equation. To make conservation manifestly Lorentz covariant, we need a fourth component that joins the three \(p^i=mu^i\). The obvious completion is
\begin{equation}
p^\mu=mu^\mu.
\label{eq:sr3-four-momentum}
\end{equation}
In \(c=1\) units,
\begin{equation}
p^\mu=(p^0,\vec p)
=
m\gamma(1,\vec v).
\end{equation}

\begin{figure}[htbp]
\centering
\includegraphics[width=0.90\textwidth]{lecture_03_figure_14.png}
\caption{Momentum conservation becomes invariant only after the spatial components are completed by a fourth component.}
\lectureframe{01:09:51}
\end{figure}

The missing component must itself be conserved. Particle mechanics already contains another conserved quantity: energy. The question is whether \(p^0=mu^0=m\gamma\) is exactly that energy in natural units. The systematic answer comes from the Hamiltonian.

\section{The Hamiltonian Is the Energy}

For any time-independent Lagrangian of this form, the Hamiltonian is
\begin{equation}
H=\sum_{i=1}^{3}\dot x^i p_i-L.
\label{eq:sr3-hamiltonian-definition}
\end{equation}
This formula is not specifically Newtonian. It is the general Legendre transform from velocities to canonical momenta and applies equally to the relativistic Lagrangian.

\begin{figure}[htbp]
\centering
\includegraphics[width=0.90\textwidth]{lecture_03_figure_15.png}
\caption{The general Hamiltonian formula \(H=\sum_i\dot x^i p_i-L\).}
\lectureframe{01:11:55}
\end{figure}

Insert \(p_i=m\dot x_i/\sqrt{1-v^2}\) and \(L=-m\sqrt{1-v^2}\):
\begin{align}
H
&=
\frac{mv^2}{\sqrt{1-v^2}}
+m\sqrt{1-v^2}\\
&=
\frac{mv^2+m(1-v^2)}{\sqrt{1-v^2}}\\
&=
\boxed{\frac{m}{\sqrt{1-v^2}}}
=m\gamma.
\label{eq:sr3-hamiltonian-result}
\end{align}
The cancellation is the key step. Comparing with \(u^0=\gamma\),
\begin{equation}
H=mu^0=p^0.
\end{equation}
Thus the fourth component is the energy:
\begin{equation}
p^\mu=(E,\vec p),
\qquad
E=H
\qquad(c=1).
\label{eq:sr3-four-momentum-natural}
\end{equation}

\begin{figure}[htbp]
\centering
\includegraphics[width=0.90\textwidth]{lecture_03_figure_16.png}
\caption{The two terms in the Hamiltonian combine over a common denominator; the \(mv^2\) contributions cancel.}
\lectureframe{01:14:28}
\end{figure}

\begin{classroomqa}
\audiencequestion{Is the formula \(H=\sum_i\dot x^i p_i-L\) only the classical Hamiltonian formula?}
\lecturerresponse{It is the general Legendre transform associated with a Lagrangian, not a special nonrelativistic assumption. The form of \(L\) decides which mechanics results. Here the relativistic square-root Lagrangian produces \(H=m/\sqrt{1-v^2}\). \lecturetimestamp{01:15:21}}
\end{classroomqa}

The invariant conservation statement is now
\begin{equation}
\sum_{\mathrm{initial}}p^\mu
=
\sum_{\mathrm{final}}p^\mu,
\label{eq:sr3-four-momentum-conservation}
\end{equation}
which contains conservation of energy and all three components of momentum in one four-vector equation.

\section{Rest Energy and the Newtonian Limit}

In natural units the exact energy is
\begin{equation}
E=\frac{m}{\sqrt{1-v^2}}.
\end{equation}
Using \eqref{eq:sr3-gamma-expansion},
\begin{equation}
E
\approx
m+\frac12mv^2+\frac38mv^4+\cdots .
\label{eq:sr3-energy-natural-expansion}
\end{equation}
The second term is familiar kinetic energy. The first is the energy that remains when the center of mass is at rest.

\begin{figure}[htbp]
\centering
\includegraphics[width=0.90\textwidth]{lecture_03_figure_17.png}
\caption{The exact energy expands as \(m+mv^2/2+\cdots\) in units \(c=1\).}
\lectureframe{01:18:02}
\end{figure}

Restore dimensions in the exact expression. The ratio inside the square root must be \(v^2/c^2\), and the numerator must have units of energy:
\begin{equation}
\boxed{
E
=
\gamma mc^2
=
\frac{mc^2}{\sqrt{1-v^2/c^2}}
}.
\label{eq:sr3-energy-exact}
\end{equation}
Its low-speed expansion is
\begin{equation}
\boxed{
E
=
mc^2+\frac12mv^2
+\frac38m\frac{v^4}{c^2}
+\cdots
}.
\label{eq:sr3-energy-dimensional-expansion}
\end{equation}
At rest,
\begin{equation}
\boxed{E_{\mathrm{rest}}=mc^2}.
\label{eq:sr3-emc2}
\end{equation}

\begin{figure}[htbp]
\centering
\includegraphics[width=0.90\textwidth]{lecture_03_figure_18.png}
\caption{Dimensional restoration turns the constant term into \(mc^2\), followed by the Newtonian kinetic energy \(mv^2/2\).}
\lectureframe{01:19:16}
\end{figure}

This applies to any closed object, not only an elementary particle. \emph{At rest} means that the center of mass is at rest. Molecules, atoms, and fields inside the object may still move and interact. Their internal kinetic energy, potential energy, and binding energy all contribute to the total rest energy and hence to the object's invariant mass.

\begin{classroomqa}
\audiencequestion{Could the conversion factor have been \(c^2\) divided by some arbitrary number?}
\lecturerresponse{Dimensions alone determine the power of \(c\), but the coefficient is fixed by the natural-unit result. The ordinary-unit formula must reduce to \(E=m\) when \(c=1\), so the restored expression is \(E=mc^2\) at rest, with coefficient one. \lecturetimestamp{01:22:08}}
\end{classroomqa}

\begin{classroomqa}
\audiencequestion{Is the leading one in the binomial expansion physically important, and what do the higher terms mean?}
\lecturerresponse{Yes. In the energy expansion the leading one becomes the rest energy \(mc^2\). The next term is Newtonian kinetic energy, and all higher terms are relativistic corrections. They are negligible for a baseball but become essential as \(v\) approaches \(c\). \lecturetimestamp{01:22:37}}
\end{classroomqa}

The exact energy and momentum diverge as \(v\to c\):
\begin{equation}
\begin{aligned}
E&=\gamma mc^2\longrightarrow\infty,\\
\vec p&=\gamma m\vec v\longrightarrow\infty
\qquad(m>0).
\end{aligned}
\end{equation}
Reaching light speed would therefore require unbounded energy and momentum. A finite force acting for finite time supplies only finite impulse, so a massive particle can approach \(c\) but cannot reach it.

\begin{classroomqa}
\audiencequestion{What now enforces the light-speed barrier?}
\lecturerresponse{The exact four-momentum does. Both energy and momentum diverge as a massive particle approaches \(c\), while four-momentum conservation constrains every interaction. The earlier normalization of four-velocity is the same statement in kinematic form. \lecturetimestamp{01:24:27}}
\end{classroomqa}

\section{The Mass of a Composite System}

Suppose two particles move with equal and opposite momenta. Their total three-momentum vanishes, so we are already in the rest frame of the combined system. The invariant mass \(M\) of that system is defined by its total energy in this frame:
\begin{equation}
Mc^2=E_{\mathrm{tot}}
\qquad
(\vec P_{\mathrm{tot}}=0).
\end{equation}
More generally,
\begin{equation}
M^2c^4
=
E_{\mathrm{tot}}^2-c^2\vec P_{\mathrm{tot}}^{\,2}.
\label{eq:sr3-composite-mass}
\end{equation}

\begin{classroomqa}
\audiencequestion{For two particles moving in opposite directions, is the mass of the combined system larger than the sum of their individual masses?}
\lecturerresponse{If they are free, yes: in the center-of-momentum frame their kinetic energy adds to the total rest energy, so \(M>m_1+m_2\). Attractive binding energy can lower the total instead, giving \(M<m_1+m_2\). The invariant definition is the total energy in the frame where total momentum vanishes. \lecturetimestamp{01:25:06}}
\end{classroomqa}

This is why internal motion and binding cannot be ignored when weighing a composite object. The invariant mass belongs to the whole closed system, not to a mere arithmetic list of constituent masses.

\begin{classroomqa}
\audiencequestion{What role does the rest-energy term play in the phase of a quantum state?}
\lecturerresponse{A state of definite energy evolves with phase \(\exp(-iEt/\hbar)\), so a particle at rest contributes \(\exp(-imc^2t/\hbar)\). A common overall phase is not observable, but relative phases between components with different energies can be. This quantum consequence is consistent with, but not needed for, the classical derivation here. \lecturetimestamp{01:26:57}}
\end{classroomqa}

\section{The Energy--Momentum Relation}

The normalization of four-velocity immediately gives the invariant norm of four-momentum. In \(c=1\) units,
\begin{equation}
(u^0)^2-\vec u^{\,2}=1.
\end{equation}

Multiplying by \(m^2\) and using \(p^\mu=mu^\mu\),
\begin{equation}
\boxed{
E^2-\vec p^{\,2}=m^2
}
\qquad(c=1).
\label{eq:sr3-mass-shell-natural}
\end{equation}
Equivalently,
\begin{equation}
E=\sqrt{\vec p^{\,2}+m^2}.
\end{equation}

\begin{figure}[htbp]
\centering
\includegraphics[width=0.90\textwidth]{lecture_03_figure_19.png}
\caption{Multiplying the four-velocity normalization by \(m^2\) yields \(E^2-\vec p^{\,2}=m^2\).}
\lectureframe{01:31:48}
\end{figure}

Restoring \(c\),
\begin{equation}
\boxed{
E^2
=
\vec p^{\,2}c^2+m^2c^4
},
\qquad
E=\sqrt{\vec p^{\,2}c^2+m^2c^4}.
\label{eq:sr3-mass-shell}
\end{equation}
This formula expresses energy as a function of momentum without referring to velocity. It is the relativistic analog of the Newtonian relation
\begin{equation}
K=\frac{\vec p^{\,2}}{2m}.
\end{equation}

\begin{figure}[htbp]
\centering
\includegraphics[width=0.90\textwidth]{lecture_03_figure_20.png}
\caption{The mass-shell relation with dimensions restored: \(E=\sqrt{\vec p^{\,2}c^2+m^2c^4}\).}
\lectureframe{01:32:38}
\end{figure}

\subsection{Massless particles}

The velocity formula \(E=\gamma mc^2\) becomes an indeterminate \(0/0\) form if \(m=0\) and \(v=c\). Velocity is therefore a poor label for distinguishing massless particles: every freely propagating massless particle in vacuum moves at \(c\), yet different photons can have different energies.

The energy--momentum relation has a smooth massless limit:
\begin{equation}
m=0
\quad\Longrightarrow\quad
\boxed{E=c|\vec p|}.
\label{eq:sr3-massless-energy}
\end{equation}
This applies to photons and, in the idealized massless theory, gravitons. It is an excellent approximation for a highly relativistic particle whenever \(mc^2\ll E\), including neutrinos at sufficiently high energy despite their small nonzero masses.

\begin{figure}[htbp]
\centering
\includegraphics[width=0.90\textwidth]{lecture_03_figure_21.png}
\caption{The massless limit of the energy--momentum relation gives \(E=c|\vec p|\).}
\lectureframe{01:33:48}
\end{figure}

\section{Positronium Decay}

Consider positronium, a bound state of an electron and a positron. Its invariant mass is slightly less than \(2m_e\): the constituents have positive internal kinetic energy but a larger negative binding energy. Numerically,
\begin{equation}
m_ec^2\approx0.511\ \mathrm{MeV},
\qquad
m_e\approx9.11\times10^{-31}\ \mathrm{kg},
\end{equation}
but the kinematics does not depend on the numerical value.

Take the two-photon decay channel in the positronium rest frame:
\begin{equation}
\mathrm{Ps}\longrightarrow\gamma+\gamma.
\end{equation}
Initially the total momentum is zero. Momentum conservation therefore requires the photons to emerge back to back with equal and opposite momenta:
\begin{equation}
\vec p_1=\vec p,
\qquad
\vec p_2=-\vec p.
\label{eq:sr3-positronium-momentum}
\end{equation}
If \(M\) is the positronium mass, energy conservation gives
\begin{equation}
Mc^2
=E_1+E_2
=c|\vec p_1|+c|\vec p_2|
=2c|\vec p|.
\end{equation}
Hence
\begin{equation}
\boxed{
|\vec p|=\frac{Mc}{2},
\qquad
E_1=E_2=\frac{Mc^2}{2}
}.
\label{eq:sr3-positronium-result}
\end{equation}

\begin{figure}[htbp]
\centering
\includegraphics[width=0.90\textwidth]{lecture_03_figure_22.png}
\caption{A positronium state at rest decays into two photons with equal and opposite momenta; \(Mc^2=2c|\vec p|\).}
\lectureframe{01:40:55}
\end{figure}

This process has no Newtonian analog in which the sum of particle masses is conserved. The incoming state has nonzero invariant mass while each outgoing photon is massless. What is conserved is total four-momentum. The stored rest energy has become propagating electromagnetic energy. With the quantum relation \(|\vec p|=h/\lambda\), the same result fixes the photon wavelength; the corresponding photon energy lies far above the optical range.

\begin{classroomqa}
\audiencequestion{How is angular momentum conserved in the decay?}
\lecturerresponse{The outgoing radiation must carry the same total angular momentum as the initial state. In the two-photon channel the photon helicities and their orbital state combine appropriately. More precisely, spin-singlet para-positronium decays predominantly into two photons, while spin-triplet ortho-positronium decays predominantly into three; the allowed channel respects all quantum numbers. \lecturetimestamp{01:42:11}}
\end{classroomqa}

\begin{classroomqa}
\audiencequestion{Why must a particle moving at light speed be massless?}
\lecturerresponse{For finite nonzero energy, the factor \(\gamma\) diverges as \(v\to c\). The invariant relation \(E^2-\vec p^{\,2}c^2=m^2c^4\) shows that a particle whose four-momentum is null has \(m=0\). Conversely, a free massless particle has a null worldline and propagates at \(c\). \lecturetimestamp{01:42:50}}
\end{classroomqa}

\begin{classroomqa}
\audiencequestion{Could the same neutral state decay into gravitons instead of photons?}
\lecturerresponse{In principle gravitational final states are allowed when the quantum numbers permit them, but the gravitational branching fraction is fantastically smaller than the electromagnetic one. The kinematic bookkeeping would still use the same conservation of four-momentum and \(E=c|\vec p|\) for each massless quantum. \lecturetimestamp{01:43:20}}
\end{classroomqa}

\subsection{One decay vertex}

A decay is local: the parent worldline ends at one spacetime event and the outgoing photon worldlines begin there. Detector interactions can be extrapolated backward to reconstruct that common vertex. If the two photons were emitted at distinct events, their back-extrapolated null trajectories would not meet at one point.

\begin{classroomqa}
\audiencequestion{Must the two photons be emitted simultaneously in every frame?}
\lecturerresponse{They originate at the same spacetime event, not at two spatially separated events whose simultaneity could be frame-dependent. All Lorentz frames agree that one event is one event. Quantum mechanics gives the parent a finite position uncertainty, so the reconstructed vertex is a distribution rather than a perfectly known classical point. \lecturetimestamp{01:44:01}}
\end{classroomqa}

\begin{classroomqa}
\audiencequestion{Why is the four-velocity norm equal to one?}
\lecturerresponse{Proper time was defined by \(d\tau^2=dt^2-d\vec x^{\,2}\) in units \(c=1\). Dividing by \(d\tau^2\) gives \((dt/d\tau)^2-(d\vec x/d\tau)^2=1\), which is precisely \((u^0)^2-\vec u^{\,2}=1\). The normalization is therefore built into the definition of proper time. \lecturetimestamp{01:45:38}}
\end{classroomqa}

\begin{classroomqa}
\audiencequestion{Does \(u^0=dt/d\tau\) give a deeper interpretation of the energy?}
\lecturerresponse{It gives the useful identity \(E=mc^2\,dt/d\tau=\gamma mc^2\): at rest the ratio of coordinate time to proper time is one, and in motion it is the Lorentz factor. Beyond encoding that velocity dependence, the identity by itself does not add a separate dynamical principle. \lecturetimestamp{01:46:36}}
\end{classroomqa}

\section{Conservation Requires a Closed System}

The statement
\begin{equation}
\sum p^\mu_{\mathrm{initial}}
=
\sum p^\mu_{\mathrm{final}}
\end{equation}
applies to an isolated system containing everything that exchanges energy and momentum. Real laboratories have stray electric and magnetic fields, ambient particles, gravity, supports, and detectors. If one leaves those degrees of freedom outside the bookkeeping, the subsystem can show a small apparent imbalance; the missing four-momentum has gone into the environment.

\begin{classroomqa}
\audiencequestion{Could the two-photon momentum sum be a tiny amount different from zero?}
\lecturerresponse{Not for the complete isolated system. In an actual experiment an unmodeled field or recoil of the apparatus can carry the tiny remainder, so the selected subsystem may not look closed. Once every interacting part is included, total four-momentum is conserved. \lecturetimestamp{01:47:51}}
\end{classroomqa}

\begin{classroomqa}
\audiencequestion{Does the uncertainty principle weaken momentum conservation?}
\lecturerresponse{No. A quantum state may have uncertain position or momentum, and sharper knowledge of one broadens the distribution of the other, but translation symmetry still gives exact conservation of total momentum for an isolated system. Quantum uncertainty concerns outcomes and distributions, not a failure of the conservation law. \lecturetimestamp{01:49:19}}
\end{classroomqa}

\section{Exact and Approximate Energy}

It is important not to confuse the exact formula with its low-speed expansion:
\begin{align}
E_{\mathrm{exact}}
&=
\frac{mc^2}{\sqrt{1-v^2/c^2}},
\label{eq:sr3-exact-repeated}\\
E_{\mathrm{slow}}
&=
mc^2+\frac12mv^2
+O\!\left(\frac{mv^4}{c^2}\right).
\label{eq:sr3-approx-repeated}
\end{align}
The first remains valid for every \(v<c\). The second is useful only when \(v/c\ll1\).

\begin{classroomqa}
\audiencequestion{Was the formula for energy only an approximation?}
\lecturerresponse{\(E=mc^2+mv^2/2+\cdots\) is the approximation. The exact result is \(E=mc^2/\sqrt{1-v^2/c^2}\). The factor \(v^2/c^2\) makes the denominator dimensionless, while the numerator \(mc^2\) gives energy units. \lecturetimestamp{01:49:56}}
\end{classroomqa}

\subsection{The spacetime picture revisited}

The late blackboard sketch returns to the geometry of the experiment. Before decay, the positronium at rest follows a vertical timelike worldline. At the decay event that worldline ends and two null photon worldlines depart in opposite spatial directions.

\begin{figure}[htbp]
\centering
\includegraphics[width=0.90\textwidth]{lecture_03_figure_23.png}
\caption{The positronium worldline terminates at one decay event, from which the two null photon worldlines emerge.}
\lectureframe{01:53:48}
\end{figure}

\begin{classroomqa}
\audiencequestion{What determines the axis along which the two photons emerge?}
\lecturerresponse{For an unpolarized state there is a quantum angular distribution, and an individual decay selects a direction randomly from that distribution. Momentum conservation fixes the photons to be back to back once one direction is realized; it does not choose that axis in advance. \lecturetimestamp{01:53:55}}
\end{classroomqa}

\section{What the Free Equations of Motion Say}

The invariant action was not chosen by guessing a force law. It was chosen as the simplest local Lorentz scalar that reproduces the Newtonian free-particle Lagrangian at low speed. Varying it gives
\begin{equation}
\frac{d}{dt}\left(\frac{\partial L}{\partial\dot x^i}\right)=0,
\end{equation}
so
\begin{equation}
\frac{dp_i}{dt}=0.
\end{equation}
For fixed nonzero mass this is equivalent to constant four-velocity:
\begin{equation}
\frac{du^\mu}{d\tau}=0.
\end{equation}
The worldline is therefore a straight line in spacetime. That single geometric statement contains both uniform speed and unchanging spatial direction.

\begin{classroomqa}
\audiencequestion{Is there a more bottom-up justification for the relativistic action, analogous to choosing the classical action so that it reproduces Newton's equations?}
\lecturerresponse{For the free particle, Lorentz invariance, locality along the worldline, and the Newtonian low-speed limit determine the simplest action up to an overall constant. Its Euler--Lagrange equation then gives constant momentum, exactly the force-free Newtonian result. Interactions provide further tests; later a charged particle will be coupled to electric and magnetic fields. \lecturetimestamp{01:54:27}}
\end{classroomqa}

\section{Two Independent Classical Limits}

The nonrelativistic limit and the classical limit of quantum mechanics are different operations. Relativistic corrections are governed by the ratio \(v/c\). Quantum effects are governed by action scales relative to \(\hbar\), together with properties such as mass, system size, and how smoothly external forces vary.

\begin{classroomqa}
\audiencequestion{Is the low-velocity limit that recovers Newtonian mechanics the same limit by which quantum mechanics becomes classical?}
\lecturerresponse{No. The relativistic limit concerns \(v/c\); the quantum-to-classical limit concerns \(\hbar\) relative to the characteristic action and often becomes effective for large masses or smoothly varying forces. A system can be quantum but nonrelativistic, relativistic but essentially classical, both, or neither. \lecturetimestamp{01:57:32}}
\end{classroomqa}

\begin{classroomqa}
\audiencequestion{Is there an incompatibility between special relativity and quantum mechanics?}
\lecturerresponse{No. Relativistic quantum theory is consistent and leads to quantum field theory. The famous unresolved tension concerns incorporating gravity into quantum mechanics, not combining special relativity with quantum principles. Gravity has not entered the present free-particle discussion. \lecturetimestamp{01:59:01}}
\end{classroomqa}

For now the force-free result is complete: a massive particle follows a straight timelike worldline, its invariant mass fixes the norm of its four-momentum, and energy and momentum are the four components of one conserved object. Introducing charge and electromagnetic fields will be the next step beyond this free motion.
